% Shared by the two standalone PDF documents. The caller defines \cacheheading.
% New displays are unnumbered to preserve the existing equation numbering.
\clearpage
\cacheheading{Persisting the FORM expansion before projection}
\code{index/form_expansion_cache.py} now owns \code{IndexFormTerm}, the exact
parser, and \code{FormExpansionCache}. For raw program text $P$, the stored value is
the parsed list of monomials before any gauge projection:
\[
 P\longmapsto \left\{\left(c,\alpha,\beta,\gamma,
        ((p,\mathbf n_p))_p\right)\right\},\qquad
 c\,t^\alpha y^\beta u^\gamma\prod_{p,j}C_p(j)^{n_{p,j}}.
\]
This notation is a direct description of the project record format. Each
\code{IndexFormTerm} has a \code{Fraction} coefficient, integer fugacity powers,
and a tuple of character indices and Adams-power tuples. The compact JSON row is
\code{["numerator", "denominator", t_power, y_power, u_power, characters]}.
\code{_encode_expansion} stores a list of these rows; \code{_decode_expansion}
restores exact arbitrary-size fractions and nested tuples. No float conversion
or representation decomposition occurs in this layer.

\begin{lstlisting}[language=SQL]
CREATE TABLE form_expansions (
    program TEXT COLLATE BINARY NOT NULL PRIMARY KEY,
    expansion_json TEXT NOT NULL
) WITHOUT ROWID;
\end{lstlisting}
The key is the complete program text, compared verbatim, not a hash. Whitespace,
numeric multiplicities, symbol names and the truncation order matter. There is
no program normalization or symbolic-multiplicity template. The default file is
\code{form_expansion_cache.db}, separate from the character database; the FORM
module does not set a separate \code{user_version}.

\textbf{Lookup.} \code{get_expansion(program)} reads and decodes a hit. A miss
runs \code{run_form}, calls \code{parse_form_output}, encodes the terms and
commits with \code{ON CONFLICT (program) DO NOTHING}. Failed execution or parsing
does not create an entry. Each thread opens its own connection, with a PID check
before reusing it, WAL, a 30-second busy timeout and up to three attempts for
busy/locked writes. FORM runs outside write transactions. Concurrent identical
misses can execute FORM more than once, but retain only one complete row.
The context manager closes the current thread's connection. These boundaries
follow the implementation and SQLite's serialized-writer model \cite{sqlite-wal,python-sqlite}.

\textbf{Index integration.} \code{calculate_index} and
\code{calculate_index_internal} accept \code{form_cache_database_path=}; the
index CLI exposes \code{--form-cache-database}. Otherwise the FORM database is
placed beside the selected character database. A hit bypasses FORM execution
and FORM-output parsing, while group-specific singlet projection still runs.
Orders below two return the vacuum without external execution or cache access.

Different theories can share this expansion when their raw programs are
identical and they use the same FORM database. Group and irrep meanings are
supplied afterward by each theory's character basis. Their final indices can
therefore differ despite a shared expansion. This reuse applies to the supported
non-$SU$ groups too; it does not require a power-sum/Schur integration backend.

\clearpage
\cacheheading{Reusing lower-order character decompositions}
Write the cached virtual character as
\[
 D_R(\mathbf n)=\sum_\lambda d_\lambda\chi_\lambda
   =\prod_{j\geq1}(\psi^j\chi_R)^{n_j},\qquad
 q(\mathbf n)=\sum_j j n_j.
\]
If $\mathbf a+\mathbf b=\mathbf n$ componentwise, multiplication in the
representation ring gives
\[
 D_R(\mathbf n)=D_R(\mathbf a)D_R(\mathbf b),\qquad
 d_\nu=\sum_{\lambda,\mu}
     a_\lambda b_\mu N_{\lambda\mu}^{\nu},
 \quad \chi_\lambda\chi_\mu=\sum_\nu N_{\lambda\mu}^{\nu}\chi_\nu.
\]
The identity follows by multiplying the defining Adams products; LiE's
\code{tensor} implements the last decomposition \cite[Sec.~4.5]{lie}.
Signed coefficients remain valid. The two exact base cases are
$D_R(\mathbf e_1)=\chi_R$ and $D_{\mathbf1}(\mathbf n)=1$;
neither needs a LiE call.

\code{CharacterDecompositionCache._decomposition_dependencies} first tries the
usual split: remove one factor at the smallest occupied Adams index. For at most
two factors, or when both usual dependencies are available, it retains that
split. Otherwise it searches lower-weighted-order keys for the \emph{same}
Cartan type and Dynkin labels, including thread-local entries and the known
first Adams operation. It considers pairs already available whose padded
exponent vectors sum to the requested vector.

Among alternative pairs, the score is the product of the two numbers of
irreducible terms, with zero cost for a zero or scalar-trivial factor. Ties use
the total number of terms and then the vectors themselves. This estimates
tensor work; it does not predict the fastest LiE call or find a globally optimal
factorization. The batch dependency planner and the actual calculation use the
same selection. Existing SQLite keys, schema and thread/process ownership are
unchanged.

For example, if $R^{\otimes2}$ and $(\psi^2R)^{\otimes2}$ are stored, their
product can be obtained with one tensor operation on those entries. Atomic
prerequisites need not be recalculated just to reconstruct a cached composite.
On-demand generation still constructs missing prerequisites when no suitable
cached split is available. It does not automatically enumerate all lower orders.

\textbf{Scope.} This optimization is separate from the singlet splitter. The
singlet path still decomposes two partial products and contracts dual labels;
the new dependency choice helps construct the required same-irrep partial
products. Intermediate full tensor decompositions can remain expensive even
when all individual Adams operations have already been cached.

\clearpage
\cacheheading{Building every product through a chosen Adams order}
\code{build_decomposition_cache(cartan_type, dynkin_labels, max_adams_order, ...)}
in \code{index/char_decomposition_cache.py} fills the existing decomposition
table for one base irrep. At weighted order $q$ it calls
\code{common.math_utils.frobenius_solve(range(1, q + 1), q)} to enumerate
\[
 \mathcal P_q=\left\{(n_1,\ldots,n_q)\in\mathbb Z_{\geq0}^q:
                 \sum_{j=1}^{q}j n_j=q\right\}.
\]
Each solution specifies one Adams product. This is also the multiplicity
description of an integer partition: $|\mathcal P_q|=p(q)$, with counts
$1,2,3,5,7,11$ at orders one through six (29 rows in total).

Orders run sequentially. The builder reads existing keys at the current order,
skips complete rows without decoding their payloads, and divides missing
solutions into batches for a persistent \code{ProcessPoolExecutor} using
\code{spawn}. Batches contain at most 64 requests, with several batches per
worker to balance uneven costs. Workers read prerequisites and calculate;
the parent commits each returned batch. It submits the next order only after
the current one is fully committed. LiE never runs in a write transaction.
Completed batches survive a later failure or interruption, so a rerun resumes
from the missing entries. Independent builders may duplicate simultaneous
misses; uniqueness and short transactions preserve complete stored rows.

Every composite at order $q$ has a split into two strictly lower orders.
Consequently both factors are already stored and at most one tensor operation
is needed for that request, with scalar/zero shortcuts as usual. For an initially
empty cache of a nontrivial irrep, one uninterrupted build needs only one
individual Adams calculation per $j=2,\ldots,M$: $M-1$ in total. This counts
logical Adams calculations, not backend retries or concurrent duplicate builds.
The builder fills same-irrep decompositions, not singlet coefficients, mixed-irrep
products or the FORM database.

\textbf{Cutoff for the full index.} A letter starts at $t^2$, and its $j$-th
Adams contribution starts no earlier than $t^{2j}$. For each base irrep in a
fixed gauge factor, its product therefore satisfies
\[
 \deg_t\ \geq\ 2\sum_j j n_j=2q
 \quad\Longrightarrow\quad
 q\leq\left\lfloor\frac{N}{2}\right\rfloor
 \quad\hbox{for an index through }t^N.
\]
This is a bound derived from the implemented single-letter functions. Use
$M=\lfloor N/2\rfloor$ for each required adjoint and matter irrep, retaining
distinct conjugate labels. It is a sufficient bound, not a claim that every
partition will occur in the projector. For product groups it holds factorwise;
one bifundamental letter contributes a character to multiple factors, so one
must not sum their separate orders and apply the same bound to that sum.
Through $t^{18}$, $M=9$ suffices. For $N<2$ there is no cache to prebuild.

\clearpage
\cacheheading{Builder API, dependencies and command-line examples}
With Sage's environment active, prepare the $A_2$ fundamental through order nine:
\begin{lstlisting}[language=bash]
sage -python -m index.char_decomposition_cache A2 \
  --dynkin-labels 1 0 --max-adams-order 9 --processes 4 \
  --cache-database /tmp/index-demo/char_decomposition_cache.db
\end{lstlisting}
Repeat for every required irrep. For $A_2$ SQCD these include $(0,1)$ and the
adjoint $(1,1)$ in addition to $(1,0)$. The group rank does not multiply the
maximum Adams order. The direct script entry point
\code{sage -python index/char_decomposition_cache.py ...} is also supported.

\begin{lstlisting}[language=Python]
from index.char_decomposition_cache import build_decomposition_cache

if __name__ == "__main__":
    totals = build_decomposition_cache(
        "A2", (1, 0), 9, processes=4,
        database_path="/tmp/index-demo/char_decomposition_cache.db",
        progress=lambda q, total, computed:
            print(q, total, computed),
    )
\end{lstlisting}
The return value is \code{dict[int, int]} mapping each order to its total product count, including reused
rows. The optional callback receives \code{(order, total, computed)} in the
parent after that order is committed. \code{processes=1} is serial; the default
is the available CPU count. The pool starts only when there is parallel work.
Use a \code{__main__} guard in scripts with multiple processes.

The maximum order and worker count must be positive integers. Cartan type,
rank and Dynkin labels are validated. \code{cache_directory=} and
\code{database_path=} are mutually exclusive. Other keywords are
\code{lie_executable=} and \code{timeout=} (600 seconds per LiE call by default;
Python accepts \code{None} for no timeout). The CLI provides
\code{--max-order} as an alias for \code{--max-adams-order}, as well as
\code{--cache-directory}, \code{--lie-executable}, and numeric \code{--timeout}.
It prints computed/reused counts and exits with status 2 on a reported failure.

The builder imports \code{frobenius_solve} lazily and therefore needs OR-Tools
only when called; ordinary index/cache lookups do not acquire that dependency.
Its enumeration and complete decompositions can grow much faster than the subset
used by one index. Prebuilding is optional, and a fully warm repeat still
enumerates solutions and checks keys even though it runs no LiE or worker pool.

Use the prepared character database and a shared FORM database for index calls:
\begin{lstlisting}[language=bash]
sage -python -m index.n2_theory_index theory.json --order 18 \
  --cache-database /tmp/index-demo/char_decomposition_cache.db \
  --form-cache-database /tmp/index-demo/form_expansion_cache.db
\end{lstlisting}
Without the last option, the FORM database defaults to that same directory.
The corresponding index Python keyword is \code{form_cache_database_path=}.
